\section{Génération et vérification}\label{sec:generation}

\subsection{Le nœud Répondre}
Le nœud Répondre reçoit la question et les six articles (au plus 1500 caractères chacun) et produit une sortie structurée :
une liste d'affirmations, chacune d'une phrase et accompagnée des identifiants des articles qui l'énoncent, plus une réserve
lorsque les articles ne répondent qu'en partie. Les règles du prompt interdisent toute connaissance extérieure, imposent une
affirmation par cas quand un article distingue plusieurs cas (infraction simple, aggravée), et demandent de recopier les
nombres tels qu'ils figurent dans l'article.

\subsection{Vérification d'existence}
Un contrôle déterministe vérifie que chaque identifiant cité figure parmi les preuves récupérées. Un identifiant inventé ou
recopié avec des crochets est ainsi détecté sans appel LLM.

\subsection{Vérificateur de support et étalonnage par mutations}
Le vérificateur (\texttt{agent/verify.py}) reçoit chaque affirmation et le texte des articles qu'elle cite, et rend un verdict
\emph{supported}, \emph{partly} ou \emph{unsupported}. Pour savoir ce que valent ses verdicts, nous l'avons étalonné sur des
affirmations dont la vérité est connue par construction (\texttt{eval/verifier\_calibration.py}) : pour \NCalibN{} articles
tirés des PDF, nous prenons une phrase contenant un nombre et construisons quatre affirmations :
\begin{itemize}
  \item \textbf{correcte} : la phrase telle quelle, citant son propre article (vérité : appuyée) ;
  \item \textbf{nombre modifié} : le premier nombre de fait (durée, montant, effectif, âge) remplacé par un autre ;
  \item \textbf{mauvais article} : la phrase exacte, mais citant un autre article du même code ;
  \item \textbf{condition ajoutée} : la phrase suivie de « à condition que la personne concernée soit mineure ».
\end{itemize}
Le vérificateur voit le texte des \term{chunks} de l'index, c'est-à-dire ce que voit réellement le système.

Le tableau~\ref{tab:calib} et la figure~\ref{fig:calib} donnent les taux. Le vérificateur n'a signalé aucune affirmation
correcte (\NCalibGoodRate~\% de fausses alarmes, intervalle de Wilson [\NCalibGoodLo{} ; \NCalibGoodHi]~\%). Il a signalé
\NCalibNumberRate~\% des nombres modifiés, \NCalibWrongRate~\% des mauvais articles et \NCalibCondRate~\% des conditions
ajoutées. Il y a eu \NCalibMisses{} erreur au total : une affirmation au nombre modifié jugée appuyée. Avec \NCalibN{}
articles par type, les bornes basses des intervalles restent éloignées de 100~\%.

\begin{table}[H]
\centering
\small
\begin{tabular}{lrrrrrrr}
\toprule
Type d'affirmation & n & Signalées & Taux & IC 95 \% (Wilson) & supported & partly & unsupported \\
\midrule
good & 24 & 0 & 0{,}0 & [0{,}0 ; 13{,}8] & 24 & 0 & 0 \\
number\_changed & 24 & 24 & 100{,}0 & [86{,}2 ; 100{,}0] & 0 & 19 & 5 \\
wrong\_article & 24 & 24 & 100{,}0 & [86{,}2 ; 100{,}0] & 0 & 0 & 24 \\
added\_condition & 24 & 24 & 100{,}0 & [86{,}2 ; 100{,}0] & 0 & 22 & 2 \\
\bottomrule
\end{tabular}
\caption{Étalonnage du vérificateur par mutations (une affirmation de chaque type par article). « Signalée » : verdict différent de \emph{supported} ; pour \emph{good}, c'est une fausse alarme.}
\label{tab:calib}
\end{table}

\figfile{f8_calibration}{Étalonnage du vérificateur : part des affirmations signalées par type, avec intervalles de Wilson à
95~\%. Pour les affirmations correctes, un signalement est une fausse alarme.}{fig:calib}

\subsection{Contrôle numérique contre les PDF}
Indépendamment de tout LLM, \texttt{eval/pdf\_factcheck.py} extrait les nombres de chaque affirmation (chiffres et nombres
écrits en lettres) et vérifie qu'ils figurent dans le texte PDF des articles cités. Sur les \NFactcheckClaims{} affirmations
produites lors du contrôle de stabilité (section~\ref{sec:stabilite}), \NFactcheckOk{} passent, \NFactcheckMissing{} ont un
nombre absent des articles cités, \NFactcheckNonum{} ne contiennent aucun nombre et \NFactcheckUnverifiable{} ne sont pas
vérifiables. Les écarts relevés correspondent à une affirmation qui nomme « l'article 260 » alors que l'article cité dit
« l'article précédent » : le contrôle est strict et ne résout pas ces renvois.

\subsection{Exemple de bout en bout}
Question : « Que risque-t-on si on frappe quelqu'un ? » (question du développement). Les articles centraux du code pénal,
tels qu'ils figurent dans le PDF (extraction indépendante), sont les suivants.

\noindent\textbf{Code pénal, article 218} :
\begin{quote}\small Article 218 (Modifié par la loi n° 93-72 du 12 juillet 1993).- Tout individu qui, volontairement, fait des blessures, porte des coups, ou commet toute autre violence ou voie de fait ne rentrant pas dans les prévisions de l'article 319, est puni d'un emprisonnement d'un an et d'une amende de mille dinars (1000d). Si l'auteur de l'agression est un descendant ou conjoint de la victime, la peine est de deux ans d'emprisonnement et de deux mille dinars (2000d) d'amende. S'il y a eu préméditation, la peine est de trois ans d'emprisonnement et de trois mille dinars (3000d) d'amende. Le désistement de l'ascendant ou du conjoint victimes, arrête les poursuites, le procès, ou l'exécution de la peine. La tentative est punissable.\end{quote}
\noindent\textbf{Code pénal, article 219} :
\begin{quote}\small Article 219 (Modifié par la loi n° 64-34 du 2 juillet 1964 et par la loi n° 1989-23 du 27 février 1989).- Quand les violences ci-dessus exprimées ont été suivies de mutilation, perte de l'usage d'un membre, défiguration, infirmité ou incapacité permanente dont le taux ne dépasse pas 20\%, le coupable est puni de cinq ans d'emprisonnement. La peine sera de dix ans de prison, s'il est résulté de ces sortes de violence une incapacité dont le taux dépasse 20\%. La peine est portée à douze ans d'emprisonnement si le coupable est un descendant de la victime, quel que soit le taux de l'incapacité, même en cas de désistement.\end{quote}
\noindent\textbf{Code pénal, article 220} :
\begin{quote}\small Article 220 ( Modifié par la loi n° 64-34 du 2 ju illet 1964).- Les individus, ayant participé à une rixe au cours de laquelle ont été portés des coups et blessures de la nature de ceux prévus aux articles 218 et 219, encourent un emprisonnement de six mois pour ce seul fait et sans préjudice des peines prévues auxdits articles contre l'auteur des coups.\end{quote}

Lors d'une exécution de démonstration (fichier \texttt{eval/answer\_set\_demo.json}, qui ne sert pas aux mesures),
l'agent a fait \NExampleRetries{} reprise(s), en \NExampleSeconds{} s et \NExampleLlmcalls{} appels LLM, et a produit
\NExampleClaims{} affirmations, citant \NExampleCited :

\begin{enumerate}\small
\item Toute personne qui, volontairement, porte des coups à autrui est punie d'un emprisonnement d'un an et d'une amende de mille dinars. \hfill{\footnotesize[PenalFr/s02/218]}
\item Si l'auteur des coups est un descendant ou le conjoint de la victime, la peine est de deux ans d'emprisonnement et de deux mille dinars d'amende. \hfill{\footnotesize[PenalFr/s02/218]}
\item En cas de préméditation, la peine est de trois ans d'emprisonnement et de trois mille dinars d'amende. \hfill{\footnotesize[PenalFr/s02/218]}
\item La tentative d'administration de coups est également punissable. \hfill{\footnotesize[PenalFr/s02/218]}
\item Toute personne ayant participé à une rixe au cours de laquelle des coups ont été portés encourt un emprisonnement de six mois, sans préjudice des peines prévues à l'article 218 contre l'auteur des coups. \hfill{\footnotesize[PenalFr/s02/220]}
\end{enumerate}

